\documentclass[a4paper,twoside]{article}

\usepackage{morewrites}

\usepackage[svgnames,dvipsnames,x11names]{xcolor}
\usepackage{graphicx}
\usepackage[english]{babel}
\usepackage[colorlinks=true, linkcolor=NavyBlue, urlcolor=NavyBlue, citecolor=NavyBlue]{hyperref}
\usepackage[a4paper]{geometry}
\usepackage{ts-fonts}
\usepackage{booktabs}
\usepackage{tabularx}
\usepackage{tcolorbox}
\usepackage{fvextra}
\usepackage{amsmath}
\usepackage{float}
\usepackage{seqsplit}
\newcolumntype{Y}{>{\centering\arraybackslash}X}
\newcolumntype{Z}{>{\raggedleft\arraybackslash}X}
\usepackage{ts-code}

\usepackage{lastpage}
\usepackage{refcount}
\makeatletter
\def\ps@plain{
\let\@oddhead\@empty
\let\@evenhead\@empty
\def\@oddfoot{
\hfil
\ifnum\getpagerefnumber{LastPage}>1\relax
\thepage
\fi
\hfil}
\let\@evenfoot\@oddfoot
}
\makeatother

\title{Diagnostics}
\author{}\date{}
\begin{document}
\maketitle
TeXSmith reports what it could not do well --- a missing image, a
cross-reference nobody publishes, a counter defined twice --- as
\textbf{diagnostics}: one line each, in the shape the \tscodeinline{tmark} toolchain prints, so a
finding reads the same whether Python or the Rust core found it.
\section{What a line looks like}
\begin{tscode}[lang=text, engine=pygments]
\begin{Verbatim}[commandchars=\\\{\},breaklines, breakanywhere, commandchars=\\\{\}]
report.md:12:5: warning ref\PYZhy{}unresolved: Counter reference \PYZsq{}@n:missing\PYZsq{} has no matching item
\end{Verbatim}
\end{tscode}
\begin{center}
\begin{tabularx}{\linewidth}{>{\raggedright\arraybackslash}X>{\raggedright\arraybackslash}X}
\toprule
\textbf{Part} & \textbf{Meaning} \\
\midrule
\tscodeinline{report.md} & the file the finding belongs to \\
\tscodeinline{12:5} & 1-based line and 1-based \textbf{byte} column of the construct (\tscodeinline{tmark check} prints the same numbers) \\
\tscodeinline{warning} & the severity \\
\tscodeinline{ref-\allowbreak{}unresolved} & a stable kebab-case code; grep for it, filter on it in CI \\
the rest & one sentence naming the construct \\
\bottomrule
\end{tabularx}
\end{center}
A finding without a position prints the file name only
(\tscodeinline{report.md: warning crossref-\allowbreak{}inventory-\allowbreak{}missing: …}), and a message that has no
file --- an engine or network failure --- prints as \tscodeinline{warning texsmith: …}.

With \tscodeinline{-\allowbreak{}v}, a suggested fix and the related locations (the other definition of
a duplicate key) follow, indented under the line.
\section{Severities}
\begin{center}
\begin{tabularx}{\linewidth}{>{\raggedright\arraybackslash}X>{\raggedright\arraybackslash}X>{\raggedright\arraybackslash}X}
\toprule
\textbf{Severity} & \textbf{Meaning} & \textbf{\tscodeinline{-\allowbreak{}-\allowbreak{}strict}} \\
\midrule
\tscodeinline{hint} & a style nit; nothing is wrong with the output & ignored \\
\tscodeinline{info} & something was done on your behalf (a root key overriding \tscodeinline{press}) & ignored \\
\tscodeinline{warning} & the output has a visible defect (\tscodeinline{[?key]}, an empty number, a missing image) & fails \\
\tscodeinline{error} & a stage could not complete; the run usually stops & fails \\
\bottomrule
\end{tabularx}
\end{center}
Warnings and errors are always shown. \tscodeinline{-\allowbreak{}q} hides hints and info lines (they
stay in the JSON dump and in the counts).

When at least one warning or error was recorded, the run ends with a summary:
\begin{tscode}[lang=text, engine=pygments]
\begin{Verbatim}[commandchars=\\\{\},breaklines, breakanywhere, commandchars=\\\{\}]
0 errors, 2 warnings
\end{Verbatim}
\end{tscode}
\section{\tscodeinline{-\allowbreak{}-\allowbreak{}strict}}
\begin{tscode}[lang=sh, engine=pygments]
\begin{Verbatim}[commandchars=\\\{\},breaklines, breakanywhere, commandchars=\\\{\}]
texsmith\PY{+w}{ }report.md\PY{+w}{ }\PYZhy{}\PYZhy{}build\PY{+w}{ }\PYZhy{}\PYZhy{}strict
\end{Verbatim}
\end{tscode}
\tscodeinline{-\allowbreak{}-\allowbreak{}strict} exits with status 1 when any warning or error was recorded. The
check runs \textbf{after the \LaTeX{} is written and before the engine runs}, so the
\tscodeinline{.tex} is there to inspect and no PDF is produced from a document with a known
hole. The same switch lives in the front matter, for documents that must never
ship with an unresolved reference:
\begin{tscode}[lang=yaml, engine=pygments]
\begin{Verbatim}[commandchars=\\\{\},breaklines, breakanywhere, commandchars=\\\{\}]
\PY{n+nt}{press}\PY{p}{:}
\PY{+w}{  }\PY{n+nt}{features}\PY{p}{:}
\PY{+w}{    }\PY{n+nt}{strict}\PY{p}{:}\PY{+w}{ }\PY{l+lScalar+lScalarPlain}{true}
\end{Verbatim}
\end{tscode}
\tscodeinline{-\allowbreak{}-\allowbreak{}strict} replaces the former \tscodeinline{PYTHONWARNINGS=error}: these findings are no
longer Python warnings and that variable no longer promotes them.
\section{\tscodeinline{-\allowbreak{}-\allowbreak{}diagnostics-\allowbreak{}json}}
\begin{tscode}[lang=sh, engine=pygments]
\begin{Verbatim}[commandchars=\\\{\},breaklines, breakanywhere, commandchars=\\\{\}]
texsmith\PY{+w}{ }report.md\PY{+w}{ }\PYZhy{}\PYZhy{}diagnostics\PYZhy{}json\PY{+w}{ }build/diagnostics.json
\end{Verbatim}
\end{tscode}
writes every recorded diagnostic as a JSON list, sorted by file and position,
for editors and CI:
\begin{tscode}[lang=json, engine=pygments]
\begin{Verbatim}[commandchars=\\\{\},breaklines, breakanywhere, commandchars=\\\{\}]
\PY{p}{[}
\PY{+w}{  }\PY{p}{\PYZob{}}
\PY{+w}{    }\PY{n+nt}{\PYZdq{}code\PYZdq{}}\PY{p}{:}\PY{+w}{ }\PY{l+s+s2}{\PYZdq{}ref\PYZhy{}unresolved\PYZdq{}}\PY{p}{,}
\PY{+w}{    }\PY{n+nt}{\PYZdq{}severity\PYZdq{}}\PY{p}{:}\PY{+w}{ }\PY{l+s+s2}{\PYZdq{}warning\PYZdq{}}\PY{p}{,}
\PY{+w}{    }\PY{n+nt}{\PYZdq{}span\PYZdq{}}\PY{p}{:}\PY{+w}{ }\PY{p}{[}\PY{l+m+mi}{0}\PY{p}{,}\PY{+w}{ }\PY{l+m+mi}{244}\PY{p}{,}\PY{+w}{ }\PY{l+m+mi}{254}\PY{p}{],}
\PY{+w}{    }\PY{n+nt}{\PYZdq{}message\PYZdq{}}\PY{p}{:}\PY{+w}{ }\PY{l+s+s2}{\PYZdq{}Counter reference \PYZsq{}@n:missing\PYZsq{} has no matching item\PYZdq{}}\PY{p}{,}
\PY{+w}{    }\PY{n+nt}{\PYZdq{}origin\PYZdq{}}\PY{p}{:}\PY{+w}{ }\PY{l+s+s2}{\PYZdq{}texsmith\PYZdq{}}\PY{p}{,}
\PY{+w}{    }\PY{n+nt}{\PYZdq{}path\PYZdq{}}\PY{p}{:}\PY{+w}{ }\PY{l+s+s2}{\PYZdq{}report.md\PYZdq{}}\PY{p}{,}
\PY{+w}{    }\PY{n+nt}{\PYZdq{}line\PYZdq{}}\PY{p}{:}\PY{+w}{ }\PY{l+m+mi}{12}\PY{p}{,}
\PY{+w}{    }\PY{n+nt}{\PYZdq{}col\PYZdq{}}\PY{p}{:}\PY{+w}{ }\PY{l+m+mi}{5}
\PY{+w}{  }\PY{p}{\PYZcb{}}
\PY{p}{]}
\end{Verbatim}
\end{tscode}
\tscodeinline{code}, \tscodeinline{severity}, \tscodeinline{span} (\tscodeinline{[file, start, end]} in bytes), \tscodeinline{message} and the
optional \tscodeinline{fix} / \tscodeinline{related} are tmark's fields; \tscodeinline{origin} says which tool found
it (\tscodeinline{texsmith} or \tscodeinline{tmark}); \tscodeinline{path}, \tscodeinline{line} and \tscodeinline{col} are the printed location,
\tscodeinline{null} when the record has none.
\section{Deprecated spellings}
A \tscodeinline{deprecated} warning names the canonical replacement:
\begin{tscode}[lang=text, engine=pygments]
\begin{Verbatim}[commandchars=\\\{\},breaklines, breakanywhere, commandchars=\\\{\}]
report.md:16:18: warning deprecated: `[\PYZca{}key]` is deprecated, write `@key`
\end{Verbatim}
\end{tscode}
These are the only findings a tool can fix for you:
\begin{tscode}[lang=sh, engine=pygments]
\begin{Verbatim}[commandchars=\\\{\},breaklines, breakanywhere, commandchars=\\\{\}]
tmark\PY{+w}{ }lint\PY{+w}{ }\PYZhy{}\PYZhy{}fix\PY{+w}{ }\PYZhy{}\PYZhy{}diff\PY{+w}{ }report.md\PY{+w}{   }\PY{c+c1}{\PYZsh{} preview}
tmark\PY{+w}{ }lint\PY{+w}{ }\PYZhy{}\PYZhy{}fix\PY{+w}{ }report.md\PY{+w}{          }\PY{c+c1}{\PYZsh{} apply in place}
\end{Verbatim}
\end{tscode}
See Migrating to TMark for the full table and its horizons.
\subsection{\tscodeinline{-\allowbreak{}-\allowbreak{}deprecated}}
Until a document has been rewritten those warnings would fail \tscodeinline{-\allowbreak{}-\allowbreak{}strict} on
their own, which is the wrong way round: a legacy spelling is a chore, not a
hole in the output. \tscodeinline{-\allowbreak{}-\allowbreak{}deprecated} sets the level the two transition codes ---
\tscodeinline{deprecated} and \tscodeinline{deprecated-\allowbreak{}frontmatter-\allowbreak{}key} --- are reported at, and nothing
else:
\begin{tscode}[lang=sh, engine=pygments]
\begin{Verbatim}[commandchars=\\\{\},breaklines, breakanywhere, commandchars=\\\{\}]
texsmith\PY{+w}{ }report.md\PY{+w}{ }\PYZhy{}\PYZhy{}build\PY{+w}{ }\PYZhy{}\PYZhy{}strict\PY{+w}{ }\PYZhy{}\PYZhy{}deprecated\PY{+w}{ }info
\end{Verbatim}
\end{tscode}
\begin{description}
\item[{\tscodeinline{warning}}] the default; they read as any other warning and fail \tscodeinline{-\allowbreak{}-\allowbreak{}strict}.
\item[{\tscodeinline{info}}] they are lowered to \tscodeinline{info} records: still printed, ignored by \tscodeinline{-\allowbreak{}-\allowbreak{}strict},
hidden by \tscodeinline{-\allowbreak{}q}.
\item[{\tscodeinline{off}}] they are dropped entirely --- absent from the terminal, from the counts and
from \tscodeinline{-\allowbreak{}-\allowbreak{}diagnostics-\allowbreak{}json}.
\end{description}
The level is applied \textbf{before} the strict check and before the JSON dump, so
what you see is what \tscodeinline{-\allowbreak{}-\allowbreak{}strict} judged. Front matter sets the same knob for a
document that must carry it into every build:
\begin{tscode}[lang=yaml, engine=pygments]
\begin{Verbatim}[commandchars=\\\{\},breaklines, breakanywhere, commandchars=\\\{\}]
\PY{n+nt}{press}\PY{p}{:}
\PY{+w}{  }\PY{n+nt}{diagnostics}\PY{p}{:}
\PY{+w}{    }\PY{n+nt}{deprecated}\PY{p}{:}\PY{+w}{ }\PY{l+lScalar+lScalarPlain}{info}
\end{Verbatim}
\end{tscode}
The CLI option wins over the front matter. Every other warning is untouched by
either.
\section{Codes}
The codes TeXSmith emits itself are listed in \tscodeinline{texsmith.diagnostics.codes}
with their default severity; a few reuse tmark's identifiers on purpose
(\tscodeinline{ref-\allowbreak{}unresolved}, \tscodeinline{label-\allowbreak{}duplicate}, \tscodeinline{crossref-\allowbreak{}inventory-\allowbreak{}missing},
\tscodeinline{crossref-\allowbreak{}inventory-\allowbreak{}stale}, \tscodeinline{include-\allowbreak{}missing}, \tscodeinline{deprecated-\allowbreak{}frontmatter-\allowbreak{}key})
because the finding is the one tmark reports once that stage runs in Rust.
The free-form \tscodeinline{texsmith} code carries the messages that predate the codes.
\end{document}
